% !TEX root = ../arbitrage.tex
\section{Related Work}
\label{sec:related-work}

\subsection{Player Boundaries and Recognition Games}

The first modeling choice in any game is the player set. \citet{aumann2025players}
argues that the entities treated as players in game-theoretic analysis need not
be merely convenient reductions of a larger game among individuals; collectives
can themselves function as units of action. The present paper studies the
opposite boundary problem. An AI system may be represented to users as if it
were a relational agent, yet it lacks the independent retaliatory channels that
would make it a strategic player in the game governing that representation.
The player set therefore remains Firm, User, and Regulator, while the system
enters as the classified object whose apparent agency affects payoffs. The
subject-retaliation parameter $\rho_S$ formalizes this boundary: when
$\rho_S>0$, the represented subject can impose off-equilibrium costs on the
arbitrageur; when $\rho_S=0$, the represented object can be priced as agentic
without being able to discipline that pricing.

\subsection{Signaling Games and Cheap Talk}

The baseline model of costless strategic communication is \citet{crawford_sobel1982}, who show that a sender with private information can transmit only coarse partitions of the type space when messages are free. The present paper extends the Crawford--Sobel framework in three directions: it adds a third strategic player (the regulator), it assigns private types to all three sides, and it treats cross-channel inconsistency---not merely imprecision---as the equilibrium pathology. \citet{spence1973} establishes that costly signals can separate types when the marginal cost of the signal is lower for the type with the higher underlying quality. The audit-trail mechanism proposed in Section~\ref{sec:mechanism} is a direct application of the Spence single-crossing condition to governance claims: high-governance firms find it cheaper to produce auditable records because they already maintain them. \citet{cho_kreps1987} supply the Intuitive Criterion as a refinement that eliminates equilibria sustained by implausible off-path beliefs; we apply it to test whether the pooling equilibrium survives scrutiny of the deflationary message. \citet{milgrom_roberts1986} show that price and advertising jointly signal product quality when false claims are costly; the present model treats marketing and liability documentation as a dual-channel analog. \citet{akerlof1970} identifies the market-for-lemons dynamic in which quality uncertainty drives adverse selection. Our pooling equilibrium is the governance analog: when governance type is unobservable and signaling is costless, low-governance firms crowd the same message space as high-governance firms, and users cannot distinguish them. \citet{myerson1979} provides the mechanism-design foundations that underwrite Section~\ref{sec:mechanism}: the problem is to make non-arbitrage incentive compatible by changing costs, rather than to rely on voluntary disclosure.

\subsection{AI Governance and Responsible AI}

Since 2018, principle-based AI ethics frameworks have proliferated across governments, industry consortia, and international organizations \citep{floridi2018ai4people,jobin2019global}. These frameworks propose norms---fairness, transparency, accountability, beneficence---but do not alter the cost structure of the signaling game. In the terms of the present model, they are cheap talk: a firm can endorse every principle without changing its governance type $\theta_F$, and endorsement carries no single-crossing cost that would separate types. The pooling equilibrium in Section~\ref{sec:bayesian} persists precisely because such declarations are costless. Algorithmic accountability proposals, including transparency mandates and explainability requirements, address \citeauthor{pasquale2015}'s concern that automated decisions operate as opaque black boxes \citep{pasquale2015}. These proposals target the opacity parameter $\theta_S$ but typically do not address cross-channel ontological inconsistency: a firm can be transparent about model architecture while maintaining incompatible ontological registers across marketing and liability documentation. \citet{hadfield2017rules} argues that regulatory institutions must be reinvented for a flat-world economy in which private governance increasingly substitutes for public law; the mechanism-design interventions in Section~\ref{sec:mechanism} operationalize that argument by specifying the parameter changes needed to shift the equilibrium. The contribution of the present paper to this literature is diagnostic: the persistence of ontological arbitrage is not a failure of moral imagination or institutional will. It is an equilibrium property of the current signal-cost structure, and it will persist under any reform that does not change the cost parameters governing the firm's message. In particular, the paper does not assume that governance rhetoric is welfare-improving; it treats such rhetoric as another cheap signal unless it changes payoffs.

\subsection{Formal Verification in Economics}

Machine verification of game-theoretic equilibria using interactive theorem provers remains rare in both the economics and AI governance literatures. The Isabelle/HOL system \citep{nipkow2002isabelle} provides a higher-order logic proof assistant with a mature library for real analysis and algebraic reasoning, but its application to economic models has been limited. Most formal results in mechanism design and signaling theory are stated and proved in the natural-language-plus-notation style standard to economic theory, with verification delegated to peer review. The present paper uses Isabelle/HOL to discharge PBE existence, Cho--Kreps survival, and the equilibrium shifts induced by each mechanism-design intervention. This serves two functions. First, it eliminates the possibility of algebraic error in the equilibrium constructions, which involve multi-player sequential rationality conditions that are tedious to verify by hand. Second, it produces machine-readable proof objects that can be inspected, extended, or challenged independently of the authors' exposition. The methodological bet is that governance-relevant claims about equilibrium behavior should meet the same verification standard applied to safety-critical software, precisely because such claims may inform cost design in regulatory systems.
